\subsection{Spectral pruning of the antiinvariant module}
\label{sec:spectral-involution}
The matrix system in Theorem~\ref{thm:involution} is useful beyond the
two-dimensional antiinvariant space. Its entry bounds impose an additional
spectral restriction that makes larger fixed-point sets accessible.

\begin{proposition}[Perron-factor exclusion]\label{prop:module-perron}
Retain the notation of Theorem~\ref{thm:involution}, and let $i=i^*$.
Suppose the nonnegative regular fusion matrix $N_i$ is irreducible.
Then no $T_i$ can have an eigenvalue whose minimal polynomial contains
$\FPdim(i)$ as a root. Moreover, if for every $x\in J$ there are
$a,y\in J$ with $N_{ax}^{y}$ odd, then $T_g=I$.
\end{proposition}
\begin{proof}
Write $B_i=(N_{ix}^{y})_{x,y\in J}$. Since $J$ is a proper subset of
the basis, the strict Perron--Frobenius inequality for a proper principal
submatrix of an irreducible nonnegative matrix gives
\[
 \rho(T_i)\leq\rho(|T_i|)\leq\rho(B_i)<\rho(N_i)=\FPdim(i).
\]
The characteristic polynomial of the integer matrix $T_i$ is rational;
thus an irreducible factor containing $\FPdim(i)$ cannot divide it.
The bounds on $T_g$ make it diagonal with entries $\pm1$.
For $a\in J$, the identity $ga=a$ gives $T_gT_a=T_a$.
A diagonal entry $-1$ would therefore force the corresponding row of
every such $T_a$ to vanish, contradicting the indicated odd entry.
\end{proof}

If $i$ generates $R\otimes\mathbb Q$, express each basis element as a
rational polynomial $p_a(i)$. The preceding proposition removes the
Perron factor from the possible minimal polynomial of $T_i$; when
$T_g=I$, every remaining factor must also divide $p_g(t)-1$.
Since $T_i\equiv B_i\pmod2$, its characteristic polynomial must reduce
to that of $B_i$ modulo two. Enumerate all products of the remaining
irreducible factors of total degree $|J|$, allowing repetitions. Each
product fixes both $\operatorname{Tr}(T_i)$ and
$\operatorname{Tr}(T_i^2)=\sum_{x,y}(T_i)_{xy}^2$.
These equalities reduce the finite entry search substantially.
Reconstruction by the $p_a$ and verification of the entire system
\eqref{eq:module-rep}--\eqref{eq:module-bound} finish the test.
The search is exhaustive because every possible module representation
has one of the enumerated rational characteristic polynomials.

\begin{proposition}\label{prop:two-four-fixed-exclusions}
The rings $\RR{8}{2945}$ and $\RR{8}{3071}$ have no categorification
as complex fusion categories.
\end{proposition}
\begin{proof}
In both rings the involution is $g=b_1$ and
$J=\{b_2,b_3,b_4,b_5\}$; all basis elements are self-dual.
All tensor identities used below are checked automatically from the
census tensors in the supplementary certificates.

For $\RR{8}{2945}$, the element $b_3$ generates the rational algebra,
its regular matrix is irreducible, and
\[
 \det(tI-N_3)=t(t-2)(t^2-23t-1)(t^2+t-1).
\]
The Perron factor is $t^2-23t-1$. The parity criterion gives $T_g=I$,
and the character with $b_3$-value zero has $g$-value $-1$, so this
character cannot occur. Thus only $t-2$ and $t^2+t-1$ remain.
The fixed block has characteristic polynomial
$t^4+t^2+1$ modulo two. Of the three degree-four products of the two
remaining factors, only $(t^2+t-1)^2$ has this reduction. Hence
\[
 \operatorname{Tr}(T_3)=-2,\qquad
 \operatorname{Tr}(T_3^2)=6,\qquad T_3^2+T_3-I=0.
\]
In the ordered basis $J$, the required odd entries of $T_3$ are the
last two diagonal entries and the symmetric pairs $(1,4),(2,3)$.
Their squared contributions already sum to six. Consequently
\[
 T_3=\begin{pmatrix}
 0&0&0&\varepsilon\\
 0&0&\delta&0\\
 0&\delta&-1&0\\
 \varepsilon&0&0&-1
 \end{pmatrix},\qquad \varepsilon,\delta\in\{1,-1\}.
\]
Rational reconstruction, reduced modulo $t^2+t-1$, gives
$T_2=-3T_3-4I$. Its first diagonal entry is $-4$, contradicting the
bound $|(T_2)_{11}|\leq N_{b_2b_2}^{b_2}=2$.

For $\RR{8}{3071}$, the regular matrix $N_3$ is irreducible and
\[
 \det(tI-N_3)=t^3(t+2)(t^2-26t+24).
\]
The last factor is the Perron factor. After removing it, the symmetric
matrix $T_3$ has eigenvalues in $\{0,-2\}$, so $T_3^2=-2T_3$.
Every diagonal entry of its fixed block is six, and therefore every
diagonal entry $d$ of $T_3$ is even. The diagonal of the matrix identity
gives
\[
 d^2+\sum_{y\ne x}(T_3)_{xy}^2=-2d.
\]
It follows that $d\in[-2,0]\cap2\mathbb Z=\{-2,0\}$, and either
choice forces all off-diagonal entries in that row to vanish. This
contradicts $N_{b_3b_2}^{b_3}=7$, which requires the corresponding
off-diagonal entry to be odd.
\end{proof}

The exact generator search checks respectively four and zero
trace-and-norm-compatible matrices for these two records. The second
search uses $b_5$ as a rational algebra generator; the displayed direct
argument uses $b_3$ to make the contradiction shorter. Both certificates
are general categorification obstructions, with no pivotal assumption.
